\textbf{Keywords:} H(s), BiQuad, Gm-C, Active-RC, OTA

\section{Introduction}\label{introduction}

The world is analog, and these days, we do most signal processing in the
digital domain. With digital signal processing we can reuse the work of
others, buy finished IPs, and probably do processing at lower cost than
for analog.

Analog signals, however, might not be suitable for conversion to
digital. A sensor might have a high, or low impedance, and have the
signal in the voltage, current, charge or other domain.

To translate the sensor signal into something that can be converted to
digital we use analog front-ends (AFE). How the AFE looks will depend on
application, but it's common to have amplification, frequency
selectivity or domain transfer, for example current to voltage.

An ADC will have a sample rate, and will alias (or fold) any signal
above half the sample rate, as such, we also must include a anti-alias
filter in AFE that reduces any signal outside the bandwidth of the ADC
as close to zero as we need.


{
\centering
\aicfig{media/l4_achai_tikz.png}
}

\emph{Figure 1: Analog signal chain from sensor through AFE
(amplification, frequency selectivity, domain transfer) to ADC}

One example of an analog frontend is the receive chain of a typical
Bluetooth radio. The signal that arrives at the antenna, or the
``sensor'', can be weak, maybe -90 dBm.

At the same time, at another frequency, there could be a unwanted
signal, or blocker, of -30 dBm

Assume for the moment we actually used an ADC at the antenna, how many
bits would we need?

Bluetooth uses Gaussian Frequency Shift Keying, which is a constant
envelope binary modulation, and it's usually sufficient with low number
of bits, assume 8-bits for the signal is more than enough.

If we assume the maximum of the ADC must fit the blocker, and the
resolution must resolve the wanted signal, we get the numbers in the
table below.

\begin{longtable}[]{@{}lll@{}}
\toprule\noalign{}
What & Power {[}dBm{]} & Voltage {[}V{]} \\
\midrule\noalign{}
\endhead
\bottomrule\noalign{}
\endlastfoot
Blocker & -30 & 7 m \\
Wanted & -90 & 7 u \\
Resolution & & Wanted/255 = 28 n \\
\end{longtable}

Then we can calculate the number of bits as

\[ \text{ ADC resolution }\Rightarrow  \ln{ \frac{7 \text{ mV}}{28 \text{ nV}} }/\ln{2} \approx 18 \text{ bits}\]

If we were to sample at 5 GHz, to ensure the bandwidth is sufficient for
a 2.480 GHz maximum frequency we can actually compute the power
consumption.

Given the Walden figure of merit of

\[FOM = \frac{P}{2^{ENOB}fs}\]

The best FOM in literature is about 1 fJ/step, so

\[P = 1\text{ fJ/step} \times 2^{18} \times 5\text{GHz} = 1.31\text{ W}\]

If we look at a typical system, like the
\href{https://www.whoop.com/eu/en/}{Whoop}. We can have a look at
teardowns, to find the battery size.

\href{https://fccid.io/2AJ2X-WS30/Internal-Photos/Internal-Photos-4265037.pdf}{Whoop
battery} is 205mAh at 3.8 V

Then we can compute the lifetime running an ADC based Bluetooth Radio

\[\text{ Hours} = \frac{ 205 \text{ mAh}}{1.31\text{ W}/3.8\text{ V}} = 0.6\text{ h}\]

I know my whoop lasts for almost a week, so it can't be what Bluetooth
ICs do.

I know a little bit about radio's, especially inside the Whoop, since it
has

\href{https://www.nordicsemi.com/Nordic-news/2022/07/the-whoop-4-uses-nordics-nrf52840-soc}{Nordic
Inside}

I can't tell you how the Nordic radio works, but I can tell you how
others usually make their radio's. The typical radio below has multiple
blocks in the AFE.


{
\centering
\aicfig{media/l4_radio_tikz.png}
}

\emph{Figure 2: Typical radio receiver frontend: LNA, complex mixer
driven by a local oscillator, anti-alias filters and ADCs for the I and
Q paths}

First is low-noise amplifier (LNA) amplifying the signal by maybe 10
times. The LNA reduces the noise impact of the following blocks. The
next is the complex mixer stage, which shifts the input signal from
radio frequency down to a low frequency, but higher than the bandwidth
of the wanted signal. Then there is a complex anti-alias filter, also
called a poly-phase filter, which rejects parts of the unwanted signals.
Lastly there is a complex ADC to convert to digital.

In digital we can further filter to select exactly the wanted signal.
Digital filters can have high stop band attenuation at a low power and
cost. There could also be digital mixers to further reduce the
frequency.

The AFE makes the system more efficient. In the 5 GHz ADC output, from
the example above, there's lot's of information that we don't use.

An AFE can reduce the system power consumption by constraining digital
information processing and conversion to the parts of the spectrum with
information of interest.

There are instances, though, where the full 2.5 GHz bandwidth has useful
information. Imagine in a cellular base station that should process
multiple cell-phones at the same time. In that case, it could make sense
with an ADC at the antenna.

What make sense depends on the application.

\section{Filters}\label{filters}\index{filters@Filters}

A filter can be fully described by the transfer function, usually
denoted by \(H(s) = \frac{\text{output}}{\text{input}}\).

Most people will today start design with a high-level simulation, in for
example Matlab, or Python. Once they know roughly the transfer function,
they will implement the actual analog circuit.

For us, as analog designers, the question becomes ``given an \(H(s)\),
how do we make an analog circuit?'' It can be shown that a combination
of 1'st and 2'nd order stages can synthesize any order filter.

Once we have the first and second order stages, we can start looking
into circuits.

\subsection{First order filter}\label{first-order-filter}\index{first order filter@First order filter}

In the book they use signal flow graphs to show how the first order
stage can be generated. By selecting the coefficients \(k_0\) ,\(k_1\)
and \(\omega_0\) we can get any first order filter, and thus match the
\(H(s)\) we want.

I would encourage you to try and derive from the signal flow graph the
\(H(s)\) and prove to your self the equation is correct.


{
\centering
\aicfig{media/l4_first_order_tikz.png}
}

\emph{Figure 3: Signal flow graph of a general first order filter with
coefficients \(k_0\), \(k_1\) and \(\omega_0\)}

Signal flow graphs are useful when dealing with linear systems.

The instructions to compute the transfer functions are

\begin{enumerate}
\def\labelenumi{\arabic{enumi}.}
\tightlist
\item
  any line with a coefficient is a multiplier
\item
  any box output is a multiplication of the coefficient and the input
\item
  any sum, well, sum all inputs
\item
  be aware of gremlins (a sudden -+ swap)
\end{enumerate}

\[ H(s) =\frac{V_o(s)}{V_i(s)}  = \frac{ k_1 s + k_0 }{s + w_o}\]

Let's call the \(1/s\) box input \(u\)

\[u = (k_0 + k_1s)V_i - \omega_0 V_o\]

\[ V_o = u/s \]

\[ u = V_o s = (k_0 + k_1s)V_i - \omega_0 V_o\]

\[ (s + \omega_0)V_o = (k_0 + k_1s)V_i\]

\[ \frac{V_o}{V_i} = \frac{k_1s + k_0}{s + \omega_0}\]

\subsection{Second order filter}\label{second-order-filter}\index{second order filter@Second order filter}

Bi-quadratic is a general purpose second order filter.

Bi-quadratic just means ``there are two quadratic equations''. Once we
match the \(k\)'s \(\omega_0\) and \(Q\) to our wanted \(H(s)\) we can
proceed with the circuit implementation.


{
\centering
\aicfig{media/l4_biquad_tikz.png}
}

\emph{Figure 4: Signal flow graph of the bi-quadratic (second order)
filter}

\[ H(s) = \frac{k_2 s^2 + k_1 s + k_0}{s^2 + \frac{\omega_0}{Q} s +
 \omega_o^2}\]

Follow exactly the same principles as for first order signal flow graph.
If you fail, and you can't find the problem with your algebra, then
maybe you need to use Maple or Mathcad.

I guess you could also spend hours training on examples to get better at
the algebra. Personally I find such tasks mind numbingly boring, and of
little value. What's important is to remember that you can always look
up the equation for a bi-quad in a book.

\subsection{How do we implement the filter
sections?}\label{how-do-we-implement-the-filter-sections}

While I'm sure you can invent new types of filters, and there probably
are advanced filters, I would say there is roughly three types. Passive
filters, that do not have gain. Active-RC filters, with OTAs, and
usually trivial to make linear. And Gm-C filters, where we use the
transconductance of a transistor and a capacitor to set the
coefficients. Gm-C are usually more power efficient than Active-RC, but
they are also more difficult to make linear.

In many AFEs, or indeed Sigma-Delta modulator loop filters, it's common
to find a first Active-RC stage, and then Gm-C for later stages.

\section{Gm-C}\label{gm-c}\index{gm-c@Gm-C}

In the figure below you can see a typical Gm-C filter and the equations
for the transfer function. One important thing to note is that this is
Gm with capital G, not the \(g_m\) that we use for small signal
analysis.

In a Gm-C filter the input and output nodes can have significant swing,
and thus cannot always be considered small signal.


{
\centering
\aicfig{media/l4_gmc_tikz.png}
}

\emph{Figure 5: Single-ended Gm-C integrator: a transconductor driving a
capacitor}

\[ V_o = \frac{I_o}{s C} = \frac{\omega_{ti}}{s} V_i \]

\[ \omega_{ti} = \frac{G_m}{C}\]

\subsection{Differential Gm-C}\label{differential-gm-c}\index{differential gm-c@Differential gm-c}

In a real IC we would almost always use differential circuit, as shown
below. The transfer function is luckily the same.


{
\centering
\aicfig{media/l4_gmc_diff_tikz.png}
}

\emph{Figure 6: Fully differential Gm-C integrator}

\[ s C V_o = G_m Vi \]

\[ H(s) = \frac{V_o}{V_i} = \frac{G_m}{sC}\]

Differential circuits are fantastic for multiple reasons, power supply
rejection ratio, noise immunity, symmetric non-linearity, but the
qualities I like the most is that the outputs can be flipped to
implement negative, or positive gain.


{
\centering
\aicfig{media/l4_gmc_diff1_tikz.png}
}

\emph{Figure 7: Differential Gm-C integrator with outputs flipped to
implement negative gain}

\[ H(s) = \frac{V_o}{V_i} = -\frac{G_m}{sC}\]

The figure below shows a implementation of a first-order Gm-C filter
that matches our signal flow graph.

I would encourage you to try and calculate the transfer function.


{
\centering
\aicfig{media/l4_gmc1st_tikz.png}
}

\emph{Figure 8: First-order Gm-C filter with transconductors \(G_{m1}\),
\(G_{m2}\) and capacitors \(C_X\) and \(C_A\)}

Given the transfer function from the signal flow graph, we see that we
can select \(C_x\), \(C_a\) and \(G_m\) to get the desired \(k\)'s and
\(\omega_0\)

\[ H(s) = \frac{ k_1 s + k_0 }{s + w_o}\]

\[ H(s) = \frac{s \frac{C_x}{C_a + C_x} + \frac{G_{m1}}{C_a + C_x}}{s +
 \frac{G_{m2}}{C_a + C_x}}\]

Below is a general purpose Gm-C bi-quadratic system.


{
\centering
\aicfig{media/l4_gmcbi_tikz.png}
}

\emph{Figure 9: General purpose differential Gm-C biquad with five
transconductors}

\[ H(s) = \frac{k_2 s^2 + k_1 s + k_0}{s^2 + \frac{\omega_0}{Q} s +
 \omega_o^2}\]

\[ H(s) = \frac{ s^2\frac{C_X}{C_X + C_B} + s\frac{G_{m5}}{C_X + C_B} + \frac{G_{m2}G_{m4}}{C_A(C_X + C_B)}}
{s^2 + s\frac{G_{m3}}{C_X + C_B} + \frac{G_{m1}G_{m2}}{C_A(C_X + C_B)} }\]

\subsection{Finding a transconductor}\label{finding-a-transconductor}\index{finding a transconductor@Finding a transconductor}

Although you can start with the Gm-C cells in the book, I would actually
choose to look at a few papers first.

The main reason is that any book is many years old. Ideas turn into
papers, papers turn into books, and by the time you read the book, then
there might be more optimal circuits for the technology you work in.

If I were to do a design in 22 nm FDSOI I would first see if someone has
already done that, and take a look at the strategy they used. If I can't
find any in 22 nm FDSOI, then I'd find a technology close to the same
supply voltage.

Start with
\href{https://ieeexplore.ieee.org/search/searchresult.jsp?queryText=Gm-C&highlight=true&returnType=SEARCH&matchPubs=true&sortType=newest&returnFacets=ALL&refinements=PublicationTitle:IEEE\%20Journal\%20of\%20Solid-State\%20Circuits}{IEEEXplore}

I could not find a 22 nm FDSOI Gm-C based circuit on the initial search.
If I was to actually make a Gm-C circuit for industry I would probably
spend a bit more time to see if any have done it, maybe expanding to
other journals or conferences.

I know of
\href{https://scholar.google.nl/citations?user=nLhKSsMAAAAJ&hl=nl}{Pieter
Harpe}, and his work is usually superb, so I would take a closer look at
A 77.3-dB SNDR 62.5-kHz Bandwidth Continuous-Time Noise-Shaping SAR ADC
With Duty-Cycled Gm-C Integrator \citeproc{ref-li23}{{[}1{]}}

And from Figure 10 a) we can see it's a similar Gm-C cell as chapter
12.5.4 in \citeproc{ref-cjm11}{{[}2{]}}.

One of my Ph.d's used the transconductor below on his master thesis
\href{https://ntnuopen.ntnu.no/ntnu-xmlui/handle/11250/2824253}{Design
Considerations for a Low-Power Control-Bounded A/D Converter}.


{
\centering
\aicfig{media/l04_ff_gm.png}
}

\emph{Figure 10: Transconductor from the master thesis: a differential
pair with common-gate cascodes to limit the Miller effect}

\section{Active-RC}\label{active-rc}\index{active-rc@Active-RC}

The Active-RC filter should be well known at this point. However, what
might be new is how the open loop gain \(A_0\) and unity gain
\(\omega_{ta}\)

\subsection{General purpose first order
filter}\label{general-purpose-first-order-filter}

Below is a general purpose first order filter and the transfer function.
I've used the conductance \(G = \frac{1}{R}\) instead of the resistance.
The reason is that it sometimes makes the equations easier to work out.

If you're stuck on calculating a transfer function, then try and switch
to conductance, and see if it resolves.

I often get my mind mixed up when calculating transfer functions. I
don't know if it's only me, but if it's you also, then don't worry, it's
not that often you have to work out transfer functions.

Once in a while, however, you will have a problem where you must
calculate the transfer function. Sometimes it's because you'll need to
understand where the poles/zeros are in a circuit, or you're trying to
come up with a clever idea, or I decide to give this exact problem on
the exam.


{
\centering
\aicfig{media/l4_activerc_first_tikz.png}
}

\emph{Figure 11: General purpose first-order active-RC filter with
parallel RC input and feedback networks}

\[ H(s) = \frac{ k_1 s + k_0 }{s + w_o}\]

\[ H(s) = \frac{  -\frac{C_1}{C_2}s -\frac{G_1}{ C_2}}{s + \frac{G_2}{ C_2}}\]

Let's work through the calculation.

\subsubsection{Step 1: Simplify}\label{step-1-simplify}

The conductance from \(V_{in}\) to virtual ground can be written as

\(G_{in} = G_1 + sC_1\)

The feedback conductance, between \(V_{out}\) and virtual ground I write
as

\(G_{fb} = G_2 + sC_2\)

\subsubsection{Step 2: Remember how an OTA
works}\label{step-2-remember-how-an-ota-works}

An ideal OTA will force its inputs to be the same. As a result, the
potential at OTA\(-\) input must be 0.

The input current must then be

\(I_{in} = G_{in} V_{in}\)

Here it's important to remember that there is no way for the input
current to enter the OTA. The OTA is high impedance. The input current
must escape through the output conductance \(G_{fb}\).

What actually happens is that the OTA will change the output voltage
\(V_{out}\) until the feedback current , \(I_{fb}\), exactly matches
\(I_{in}\). That's the only way to maintain the virtual ground at 0 V.
If the currents do not match, the voltage at virtual ground cannot
continue to be 0 V, the voltage must change.

\subsubsection{Step 3: Rant a bit}\label{step-3-rant-a-bit}

The previous paragraph should trigger your spidy sense. Words like
``exactly matches'' don't exist in the real world. As such, how closely
the currents match must affect the transfer function. The open loop gain
\(A_0\) of the OTA matters. How fast the OTA can react to a change in
voltage on the virtual ground, approximated by the unity-gain frequency
\(\omega_{ta}\) (the frequency where the gain of the OTA equals 1, or 0
dB), matters. The input impedance of the OTA, whether the gate leakage
of the input differential pair due to quantum tunneling, or the
capacitance of the input differential pair, matters. How much current
the OTA can deliver (set by slew rate), matters.

Active-RC filter design is ``How do I design my OTA so it's good enough
for the filter''. That's also why, for integrated circuits, you will not
have a library of OTAs that you just plug in, and they work.

I would be very suspicious of working anywhere that had an OTA library I
was supposed to use for integrated filter design. I'm not saying it's
impossible that some company actually has an OTA library, but I think
it's a bad strategy. First of all, if an OTA is generic enough to be
used ``everywhere'', then the OTA is likely using too much power,
consumes too much area, and is too complex. And the company runs the
risk that the designer has not really checked that the OTA works
properly in the filter because ``Someone else designed the OTA, I just
used in my design''.

But, for now, to make our lives simpler, we assume the OTA is ideal.
That makes the equations pretty, and we know what we should get if the
OTA actually was ideal.

\subsubsection{Step 4: Do the algebra}\label{step-4-do-the-algebra}

The current flowing from \(V_{out}\) to virtual ground is

\[I_{out}= G_{out}V_{out}\]

The sum of currents into the virtual ground must be zero

\[ I_{in} + I_{out} = 0\]

Insert, and do the algebra

\[ G_{in}V_{in} + G_{out}V_{out} = 0\]

\[ \Rightarrow - G_{in} V_{in} = G_{out} V_{out}\]

\[ \frac{V_{out}}{V_{in}} = - \frac{G_{in}}{G_{out}}\]

\[ = - \frac{ G_1 + s C_1 }{G_2 + sC_2}\]

\[ = \frac{ -s \frac{C_1}{C_2} - \frac{G_1}{C_2} }{s + \frac{G_2}{C_2}}\]

\subsection{General purpose biquad}\label{general-purpose-biquad}\index{general purpose biquad@General purpose biquad}

A general bi-quadratic active-RC filter is shown below. These kind of
general purpose filter sections are quite useful.

Imagine you wanted to make a filter, any filter. You'd decompose into
first and second order sections, and then you'd try and match the
transfer functions to the general equations.

The
\href{https://wulffern.github.io/aic2026/assets/examples/biquad.html}{interactive
biquad} is worth a few minutes here: the low-pass, band-pass, high-pass,
notch and all-pass responses all share one denominator, and switching
between them moves the zeros while the poles stay exactly where they
are.


{
\centering
\aicfig{media/l4_activebiquad_tikz.png}
}

\emph{Figure 12: General purpose active-RC biquad built from two OTAs}

\[ H(s) = \frac{k_2 s^2 + k_1 s + k_0}{s^2 + \frac{\omega_0}{Q} s +
 \omega_o^2}\]

\[H(s) = \frac{\left[ \frac{C_1}{C_B}s^2 + \frac{G_2}{C_B}s + (\frac{G_1G_3}{C_A C_B})\right]}{\left[ s^2  + \frac{G_5}{C_B}s + \frac{G_3 G_4}{C_A C_B}\right]}\]

\section{The OTA is not ideal}\label{the-ota-is-not-ideal}\index{the ota is not ideal@The OTA is not ideal}


{
\centering
\aicfig{media/l4_activerc_tikz.png}
}

\emph{Figure 13: Active-RC integrator: OTA with resistor input and
capacitor feedback}

\[ H(s) \approx \frac{A_0}{(1 + s A_o R C)(1 + \frac{s}{w_{ta}})}\]

where \[A_0\] is the gain of the amplifier, and \[\omega_{ta}\] is the
unity-gain frequency.

At frequencies above \(\frac{1}{A_0RC}\) and below \(w_{ta}\) the
circuit above is a good approximation of an ideal integrator.

See page 511 in \citeproc{ref-cjm11}{{[}2{]}} (chapter 5.8.1)

\section{Example circuit}\label{example-circuit}\index{example circuit@Example circuit}

One place where both active-RC and Gm-C filters find a home are
continuous time sigma-delta modulators. More on SD later, for now, just
know that SD is a combination of high-gain, filtering, simple ADCs and
simple DACs to make high resolution analog-to-digital converters.

One such an example is

\href{https://ieeexplore.ieee.org/stamp/stamp.jsp?tp=&arnumber=4381437}{A
56 mW Continuous-Time Quadrature Cascaded Sigma-Delta Modulator With 77
dB DR in a Near Zero-IF 20 MHz Band}

Below we see the actual circuit. It may look complex, and it is.

Not just ``complex'' as in complicated circuit, it's also ``complex'' as
in ``complex numbers''.

We can see there are two paths ``i'' and ``q'', for ``in-phase'' and
``quadrature-phase''. The fantastic thing about complex ADCs is that we
can have a-symmetric frequency response around 0 Hz.

It will be tricky understanding circuits like this in the beginning, but
know that it is possible, and it does get easier to understand.

With a complex ADC like this, the first thing to understand is the rough
structure.

There are two paths, each path contains 2 ADCs connected in series
(Multi-stage Noise-Shaping or MASH). Understanding everything at once
does not make sense.

Start with ``Vpi'' and ``Vmi'', make it into a single path (set Rfb1 and
Rfb2 to infinite), ignore what happens after R3 and DAC2i.

Now we have a continuous time sigma delta with two stages. First stage
is a integrator (R1 and C1), and second stage is a filter (Cff1, R2 and
C2). The amplified and filtered signal is sampled by the ADC1i and fed
back to the input DAC1i.

It's possible to show that if the gain from \(V(Vpi,Vmi)\) to ADC1i
input is large, then \(Y1i = V(Vpi,Vmi)\) at low frequencies.

Then put \(R_{fb1}\) and \(R_{fb2}\) back. Each one takes the
\emph{first integrator output} of one path and drives it into the
\emph{summing node} of the other - the same virtual ground the input
resistor \(R_1\) and the feedback DAC already inject into. That is the
whole of the cross-coupling, and it is what turns two real modulators
into one complex one.

The figure draws the first stage of the cascade only. What leaves
through \(R_3\) and \(DAC2\) is the quantization error of this stage,
handed to a second modulator that digitizes it; the digital outputs are
then recombined so the first stage's error cancels. That is what MASH
means, and it is a separate story from the complex part.


{
\centering
\aicfig{media/qt_sd_tikz.png}
}

\emph{Figure 14: One stage of a continuous-time quadrature sigma-delta
modulator, both paths drawn: two ordinary real modulators - R1/C1
integrating, Cff1/R2/C2 filtering, ADC1 sampling, DAC1 feeding back, R3
and DAC2 handing off to the second stage - joined by the red
cross-coupling that makes the pair complex: each \(R_{fb}\) carries one
path's first integrator output into the other path's summing node.
Inspired by Breems et al. \citeproc{ref-breems07}{{[}3{]}}, which is
fully differential, cascaded, and has a tuning resistor on everything}

What does the cross-coupling actually buy? Figure 15 answers it by
simulation. The same first-order loop is run twice: once real, so its
coefficients are real and its noise notch must sit symmetrically about
zero, and once complex, with the integrator pole rotated to
\(e^{j\omega_0}\) so the noise transfer function has its zero at
\(+f_0\) and nowhere else. The complex quantizer is just the two real
quantizers, one per path - which is what the i and q ADCs are.

Read the two spectra over the \emph{whole} sample rate, not half of it.
The real loop is quiet at DC and equally noisy at \(\pm f_0\): a real
signal's spectrum is conjugate symmetric and cannot tell \(+f\) from
\(-f\). The complex loop is quiet at \(+f_0\) and noisy at \(-f_0\), so
a wanted channel just above the local oscillator gets the quiet side
while its image gets the noise. That asymmetry is the entire reason a
near zero-IF receiver bothers with quadrature.

Both simulations are driven by the same input, a complex tone just above
\(f_0\), and you can find it in both spectra as the tall spike. In the
real loop it stands out in the middle of the noise, where the loop has
done nothing for it. In the complex loop it stands in the notch. The
signal did not move; the quiet part of the spectrum did.


{
\centering
\aicfig{media/l04_complex_psd_tikz.png}
}

\emph{Figure 15: Simulated output spectra of a first-order sigma-delta
loop over the full sample rate. (a) A real loop: the noise notch is
symmetric about zero. (b) The same loop made complex: the notch moves to
\(+f_0\) alone, leaving \(-f_0\) noisy - which is what lets a near
zero-IF receiver keep the wanted side and throw the image away. The same
input tone is present in both, and only in (b) does it land in the
notch}

\section{My favorite OTA}\label{my-favorite-ota}\index{my favorite ota@My favorite OTA}

Over the years I've developed a love for the current mirror OTA. A
single stage, with load compensation, and an adaptable range of DC
gains.

Sometimes simple current mirrors are sufficient, sometimes cascoded, or
even active cascodes are necessary.

Below is the differential current mirror OTA.


{
\centering
\aicfig{media/l04_ota_diff_tikz.png}
}

\emph{Figure 16: Differential current mirror OTA}

In a differential OTA we need to control the output common mode. In
order to control the common mode, we must sense the common mode.

Below is a circuit I often use to sense the common mode. Ideally the
source followers would be native transistors, but those are not always
available.

The reference for the common mode can be from a bandgap, or in the case
below, VDD/2.


{
\centering
\aicfig{media/l04_ota_vsens_tikz.png}
}

\emph{Figure 17: Common mode sense circuit with source followers and the
resistor divider generating the reference \(V_{CREF}\)}

Once we have both the sensed common mode, and the common mode reference,
we can use another OTA to control the common mode.

The nice thing about the circuit below is that the common mode feedback
loop has the same dominant pole as the differential loop.


{
\centering
\aicfig{media/l04_ota_vcmfb_tikz.png}
}

\emph{Figure 18: Common mode feedback OTA comparing the sensed
\(V_{COUT}\) to the reference \(V_{CREF}\)}

You can find the schematic for the OTA at

\href{https://github.com/wulffern/cnr_ota_sky130nm}{CNR\_OTA\_SKY130NM}


{
\centering
\aicfig{media/l04_ota_sch.png}
}

\emph{Figure 19: CNR\_OTA schematic in SKY130: bias, differential OTA,
common mode sense (VCM) and common mode feedback OTA}

\section{Summary}\label{summary}

The one-page version of this chapter:

\begin{itemize}
\tightlist
\item
  The frontend's job is to hand the ADC a signal it can afford to
  convert: gain where it is cheap, filtering before sampling folds noise
  in
\item
  First order sections come from one gm and one C; biquads stack them
  with feedback and give Q
\item
  Gm-C is fast and open loop, active-RC is linear and closed loop - the
  OTA pays either way
\item
  A real OTA's finite gain and bandwidth move the filter poles; design
  the OTA from the filter's error budget
\item
  Fully differential filters double swing and cancel even harmonics, and
  carry the CMFB tax
\item
  A complex (quadrature) filter is two real paths plus cross-coupling:
  the response stops being symmetric about zero, which a low-IF radio
  needs
\end{itemize}

\section{Would you like to know
more?}\label{would-you-like-to-know-more}

A 77.3-dB SNDR 62.5-kHz Bandwidth Continuous-Time Noise-Shaping SAR ADC
With Duty-Cycled Gm-C Integrator \citeproc{ref-li23}{{[}1{]}}

\href{https://ntnuopen.ntnu.no/ntnu-xmlui/handle/11250/2824253}{Design
Considerations for a Low-Power Control-Bounded A/D Converter}

A 56 mW Continuous-Time Quadrature Cascaded Sigma-Delta Modulator With
77 dB DR in a Near Zero-IF 20 MHz Band \citeproc{ref-breems07}{{[}3{]}}

Complex signal processing is not - complex
\citeproc{ref-martin03}{{[}4{]}}

\phantomsection\label{refs}
\begin{CSLReferences}{0}{0}
\bibitem[\citeproctext]{ref-li23}
\CSLLeftMargin{{[}1{]} }%
\CSLRightInline{H. Li, Y. Shen, E. Cantatore, and P. Harpe, {``A 77.3-dB
SNDR 62.5-kHz bandwidth continuous-time noise-shaping SAR ADC with
duty-cycled gm-c integrator,''} \emph{IEEE Journal of Solid-State
Circuits}, vol. 58, no. 4, pp. 939--948, 2023, doi:
\href{https://doi.org/10.1109/JSSC.2022.3227678}{10.1109/JSSC.2022.3227678}.
}

\bibitem[\citeproctext]{ref-cjm11}
\CSLLeftMargin{{[}2{]} }%
\CSLRightInline{T. C. Carusone, D. Johns, and K. Martin, \emph{Analog
integrated circuit design}. Wiley, 2011 {[}Online{]}. Available:
\url{https://books.google.no/books?id=1OIJZzLvVhcC}}

\bibitem[\citeproctext]{ref-breems07}
\CSLLeftMargin{{[}3{]} }%
\CSLRightInline{L. J. Breems, R. Rutten, R. H. M. van Veldhoven, and G.
van der Weide, {``A 56 mW continuous-time quadrature cascaded
\(\Sigma\Delta\) modulator with 77 dB DR in a near zero-IF 20 MHz
band,''} \emph{IEEE Journal of Solid-State Circuits}, vol. 42, no. 12,
pp. 2696--2705, 2007, doi:
\href{https://doi.org/10.1109/JSSC.2007.908765}{10.1109/JSSC.2007.908765}.
}

\bibitem[\citeproctext]{ref-martin03}
\CSLLeftMargin{{[}4{]} }%
\CSLRightInline{K. Martin, {``Complex signal processing is not -
complex,''} in \emph{ESSCIRC 2004 - 29th european solid-state circuits
conference (IEEE cat. no.03EX705)}, 2003, pp. 3--14, doi:
\href{https://doi.org/10.1109/ESSCIRC.2003.1257061}{10.1109/ESSCIRC.2003.1257061}.
}

\end{CSLReferences}
